\documentclass[11pt,a4paper]{article}
\usepackage[margin=27mm]{geometry}
\usepackage{fontspec,unicode-math}
\setmainfont{STIXGeneral}\setmathfont{STIX Two Math}
\usepackage{amsmath,amsthm,hyperref}
\title{Independent review of the Topic F edition}
\author{Independent project review}
\date{11 October 2026}
\begin{document}\maketitle
\noindent\textbf{Decision.} The frozen common paper and application recover the already admitted corrected principal results of Hou, Feng, Leemans and Qu, \emph{String C-groups of order $4p^m$}, arXiv:2407.10388v1. No mathematical correction is requested. This accepts the stated twenty-five-record boundary, including retained and excluded records; it does not mean that twenty-five independent new results, every assertion of the source, or a new conjecture have been proved. The author manifest is bound in the accompanying JSON record. The reviewer did not author this canonical edition and has made no change to it.

\section*{Mathematical reading}
The review read both complete papers and all eight TeX sections, the entire twelve-page source, the exact current source map, all three mathematical programs and the five Lean proofs. The following checks concern the quantified arguments, independently of successful sample computations.

The rank-three intersection criterion checks all subgroup pairs, not only the two displayed maximal parabolics. The remaining two pairs of rank-two subgroups are controlled by the four-element endpoint group. Distinct nonidentity generators exclude the unwanted intersections. The resulting product formula gives the exact inequality $2k_1k_2\leq |G|$. The prime-index argument uses the actual relative subgroup index; finite ambient order is unnecessary there, but noncontainment is essential.

The normal Sylow proof handles the only exceptional count, $p=3$ and four Sylow subgroups. The conjugation kernel lies in the normalizer of order $3^m$; an image of order twelve in $S_4$ is $A_4$, which cannot be generated by its involutions. The endpoint Klein four group then really is a complement. Nondegeneracy makes both adjacent rotations have nontrivial $p$-parts. The quotient excludes two odd labels, while the intersection property excludes a cyclic Sylow subgroup.

The inverted orbit lemma makes $\langle u,u^k\rangle$ invariant under the complement and uses generation to identify it with the whole kernel. It does not claim that inversion is an automorphism of a nonabelian group. In the abelian case the four projectors are genuine orthogonal idempotents because multiplication by four is invertible. Their algebraic calculation also contains the ordinary decomposition for an involution used in source Lemma 2.4; that elementary part does not use the later orbit-generation hypothesis. The orbit hypothesis then makes the two surviving components cyclic and nontrivial.

In the abelian application, the third possible quotient involution gives a product bound twice the group order and is excluded. For the remaining marking the inequality forces the larger exponent onto the correct character component. The converse construction proves generation and the intersection property explicitly. Marked uniqueness among these models follows from their chosen cyclic generators and characters. The further implication from an arbitrary tight marking uses the stated Cunningham--Pellicer odd/even classification, with even parameter four excluded. This is a retained theorem, not an inference from the finite examples.

For the cyclotomic calculation, cancellation in the integral domain proves every successive ideal index. The ratios of the conjugate uniformizers are integral units by their norm and integral multiplication matrices; hence $pO=\pi^{p-1}O$. The diagonal lattice has the correct inclusion and index, which proves equality and every individual coordinate modulus, including modulus one. This step establishes more than an abstract group order. Complex conjugation preserves the ideal, and multiplication by $\zeta$ has exact order $p$ when $n\geq2$.

The semidirect-product convention is checked against the displayed pair law: $v^\beta=\zeta v$. The two proposed maps preserve multiplication and generate an actual Klein four group of automorphisms. The norm sum proves that every element outside the abelian subgroup has order $p$. The marked elements are distinct nonidentity involutions, their adjacent orders are exactly $p,2p$, and they generate the entire group. The prime-index lemma then gives the full intersection property through the already proved rank-three criterion.

The presentation proof establishes a surjection onto the ring model and an independent upper bound from its finite abelian normal subgroup. Both bounds are $p^m$, so the map is an isomorphism on the actual named generators. Conjugation and binomial expansion recover both clauses of source Lemma 5.2, and the ring automorphisms recover Lemma 5.3. The large coefficient table in the source is not needed in this deduction. The lower central series $\gamma_j(P)=\pi^{j-1}A$ proves the asserted maximal class, and the stated small-exponent examples correctly show why the motivating Hughes condition is unnecessary.

\section*{Source and attribution}
The literal marking on source page 2 has square $x^2$, so it cannot be recovered as an involution. The source's own character case on page 6 uses the corrected endpoint; the canonical application states both this change and the necessary exponent assignment. Source pages 6--8 were inspected visually for the presentation, marking and auxiliary formulas. The full twelve-page extracted text was read, including the source's earlier classification inputs and its excluded introductory census.

The principal structural, abelian and existence results retain their original authorship. The orbit argument follows the source's proof, the character splitting is classical, and the cyclotomic model is offered as a reusable proof organization without a historical priority claim. The present review accepts the retained Cunningham--Pellicer input on the basis of the previously admitted exact theorem record and the source's explicit invocation. Its full PDF was not reacquired or independently read in this turn; its local hash and complete proof cost remain unknown. The usual regular-polyhedron/string-C-group correspondence is likewise retained. Neither the survey census nor a classification of Hughes groups is included.

\section*{Computation and formal scope}
A fresh portable replay reproduced seven cyclotomic and four abelian marked examples, and five literal source presentations. Each marking was tested on all sixty-four subset intersections. The incorrect marking, corrupted relation and altered generator orders were rejected, and a composite-index counterexample was verified. These three complete programs validate the written proofs; they are not prerequisites for their all-parameter conclusions. Both complete historical drivers remain separately visible. One is staged and executed by the canonical runner; the other retains its original package-layout requirement and is archival.

The actual Lean statements are two subgroup intersection theorems and three group identities. The first reflection identity does not establish that a source marker is nonidentity, and the adjacent-product identity does not establish its exact order. Those facts are proved in the paper. The literal bad-marking identity proves the square formula, while its nontriviality at the asserted parameters is written. Fresh compilation reports only the declared standard axioms. The same five declarations are compiled through two adapters but are counted only once. No complete Sylow, character, lattice, construction or classification formalization is asserted.

A separate review program works over the integral $1,\zeta,\ldots,\zeta^{p-2}$ basis, imports no author code, and compares Hermite normal forms of the actual ideals with the claimed coordinate lattices. All forty cases $p=3,5,7,11$ and $1\leq n\leq10$ pass, as do the exact conjugation identities; deleting the constant cyclotomic coefficient is rejected. This is independent bounded evidence, not the proof of the uniform lattice theorem.

All thirteen current PDF pages were inspected in rendered contact sheets; no clipping or missing mathematical glyph was seen. The three fresh renders have the same text page by page as the frozen PDFs, and the complete code text is verified in the appendix. Eight integrity controls passed. These checks supplement the full mathematical reading and do not replace it. No source credit, question credit or compression ratio is added by this review.
\end{document}
